\problemname{Fluortanten}
Björn and $n-1$ other people are standing in a queue to see the fluoride lady. Different people find it differently scary to see the fluoride lady. The people are numbered from $1$ to $n$, and person $i$ stands at position $i$ in the queue. Person $i$ also has a value $a_i$, which shows how reluctant the person is to see the fluoride lady. Person $i$'s happiness about their position in the queue is $i \cdot a_i$. Some people may have negative $a_i$,
which means they actually want to see the fluoride lady and are thus sad about having to wait.

Björn is the only person who is completely indifferent to seeing the fluoride lady, that is, he is the only person with $a_i = 0$. Moreover, he is very kind-hearted, so he decides to leave the queue and then re-enter it at some position so that the total happiness of everyone in the queue is maximized. Write a program that, given the values $a_i$ for all people, computes the maximum sum of happiness in the queue if Björn places himself at an optimal position.

\section*{Input}
The first line contains an integer $n$ ($1 \le n \le 10^6$), the number of people in the queue. The next line contains $n$ integers, where the $i$th number, $-1000 \le a_i \le 1000$, describes the happiness of the $i$th person in the queue.

\section*{Output}
Print a single integer: the maximum total happiness in the queue.

\section*{Scoring}
Your solution will be tested on a set of test groups.
To earn points for a group, you must pass all test cases in that group.

\noindent
\begin{tabular}{| l | l | p{12cm} |}
  \hline
  Group & Points & Constraints \\ \hline
    $1$   & $32$     & $1 \leq n \leq 1000$ \\ \hline
    $2$   & $12$     & The sequence $a_1 \dots a_n$ is non-decreasing. \\ \hline
    $3$   & $11$     & The sequence $a_1 \dots a_n$ is non-increasing. \\ \hline
    $4$   & $45$     & No additional constraints. \\ \hline
\end{tabular}

\section*{Explanation of Sample 1}

Björn places himself last in the queue. The sum of happiness becomes $1 \cdot 1 + 2\cdot(-2) + 3 \cdot 0 = -3$.

If Björn had instead placed himself first in the queue, the sum of happiness would have been $1 \cdot 0 + 2 \cdot 1 + 3 \cdot (-2) = -4$, and in the middle it would have been $1 \cdot 1 + 2 \cdot 0 + 3 \cdot (-2) = -5$. Both of these are worse alternatives.
